\documentclass[12pt]{article}
\usepackage[utf8]{inputenc}
\usepackage{graphicx}
\usepackage{amsmath}
\usepackage{amssymb}
\usepackage{mathrsfs}
\usepackage{pgfornament}
\usepackage[many]{tcolorbox}
\usepackage[margin = 0.5in]{geometry}
\usepackage{collectbox}
\usepackage{mdframed}
\usepackage{xcolor}

\newcommand\textbox[1]{%
  \parbox{.333\textwidth}{#1}%
}

\linespread{2}

\makeatletter
\newcommand{\mybox}{%
    \collectbox{%
        \setlength{\fboxsep}{2pt}%
        \fbox{\BOXCONTENT}%
    }%
}
\makeatother

\usepackage{listings}
\usepackage{color}

\definecolor{dkgreen}{rgb}{0,0.6,0}
\definecolor{gray}{rgb}{0.5,0.5,0.5}
\definecolor{mauve}{rgb}{0.58,0,0.82}

\lstset{frame=tb,
  language=R,
  aboveskip=3mm,
  belowskip=3mm,
  showstringspaces=false,
  columns=flexible,
  basicstyle={\small\ttfamily},
  numbers=none,
  numberstyle=\tiny\color{gray},
  keywordstyle=\color{blue},
  commentstyle=\color{dkgreen},
  stringstyle=\color{mauve},
  breaklines=true,
  breakatwhitespace=true,
  tabsize=3
}
\title{MAT 523 HW 4}
\begin{document}

\begin{center}
\begin{large}
\textbf{MAT 523: Functions of a Real Variable I \\
Homework IV \\}
\end{large}
Nolan R. H. Gagnon \\
\end{center}

\noindent \textbf{(1.)}

\vspace{3mm}

\noindent\mybox{\colorbox{lightgray}{\textbf{Theorem}}}\hspace{3mm}\hrulefill

\noindent \textbf{(a)} Every real sequence has a monotone subsequence. \textbf{(b)} Use this to provide another proof of the Bolzano-Weierstrass Theorem. 

\vspace{5mm}

\noindent\mybox{\colorbox{lightgray}{\textbf{Proof}}}\hspace{3mm}\hrulefill

\noindent Suppose $E\neq\mathbb{R}$. Then $E$ must be bounded above, bounded below, or both. Without loss of generality, suppose $E$ is bounded above. Then, since $E$ is nonempty, it must have a supremum $\ell = \sup E$. Now, it is either the case that $\ell\in E$ or $\ell\not\in E$.  Suppose, first, that $\ell\in E$. Since $E$ is open, there must be an open set $(a,b)\subseteq E$ such that $\ell\in(a,b).$ But for $\ell$ to be in $(a,b),$ $\ell$ must be less than $b$. But $b\in E$, which contradicts the fact that $\ell$ is an upper bound for $E$.  Next, suppose that $\ell\not\in E$. Let $(c,d)$ be an open set containing $\ell$. If $E\cap(c,d)=\emptyset$, then all elements of $E$ must be less than or equal to $c$. But then $c$ is an upper bound for $E$ that is less than $\ell$. This is a contradiction, so $E\cap(c,d)\neq\emtpyset$. Thus, $\ell$ is a closure point of $E$. But $\ell\not\in E$, meaning $E$ is not closed. This is a contradiction.  Therefore, $E = \mathbb{R}.$

\hfill$\blacksquare$

\pagebreak

\noindent \textbf{(3.)} 

\vspace{3mm}

\noindent\mybox{\colorbox{lightgray}{\textbf{Theorem}}}\hspace{3mm}\hrulefill 

\noindent A compact set in $\mathbb{R}$ is closed and bounded. 

\vspace{5mm}

\noindent\mybox{\colorbox{lightgray}{\textbf{Proof}}}\hspace{3mm}\hrulefill

\noindent Let $E\subseteq\mathbb{R}$ be a compact set. If $E$ is empty, then by Proposition 12, $E$ is closed. It can also be argued vacuously that $E$ is bounded above and below. Now suppose $E$ is non-empty. Let $\mathscr{C}$ be an open cover of $E$. Since $E$ is compact, $\mathscr{C}$ has a finite subcover of $E$, call it $F$. Since $F$ contains finitely many open sets, we can write them all down $$(a_1, b_1), (a_2, b_2), \ldots (a_n, b_n)$$ so that $a_1 < a_2 < \ldots < a_n.$ Since $E\subseteq (a_1,b_1)\cup (a_2, b_2) \ldots \cup (a_n, b_n)$ (by definition of an open cover), it is clearly bounded below by $a_1$. We can also write down the open sets in $F$ $$(c_1, d_1), (c_2, d_2) \ldots (c_n, d_n)$$ so that $d_1 < d_2 < \ldots < d_n$.  (Obviously, the elements of this list are identical to the elements of the list above. We simply use different notation to allow for the possibility of a different ordering.) Again, since $E\subseteq (c_1, d_1) \cup (c_2, d_2) \ldots \cup (c_n, d_n)$, it must be bounded above by $d_n$.  Since $E$ is bounded above and below, $E$ is bounded. 

\vspace{3mm}

\noindent Next we show that $E$ is closed. If not, then there exists an $x\in\overline{E}$ such that $x\not\in E$. Since $x$ is a closure point of $E$, there exists an $(a,b)$ such that $x\in(a,b)$ and $E\cap(a,b)\neq\emptyset$. Thus, there is a $y\in E$ that is also in $(a,b)$. Let $\mathscr{C}$ be an open cover of $E$. Since $E$ is compact, there is a finite subcover of $E$, call it $\mathscr{F}$.

\pagebreak

\noindent \textbf{(4.)}

\vspace{3mm}

\noindent\mybox{\colorbox{lightgray}{\textbf{Problem}}}\hspace{3mm}\hrulefill

\noindent Show that the conclusions of the Nested Set Theorem can fail to hold if \textbf{(a)} instead of compact sets, we take sets that are merely closed. Likewise, show that it can fail if \textbf{(b)} the sets are taken to be open and bounded.

\vspace{5mm}

\noindent\mybox{\colorbox{lightgray}{\textbf{Solution}}}\hspace{3mm}\hrulefill

\noindent \textbf{(a)} For the first case, consider the sets $$F_n = [n,\infty),$$ where $n\in\mathbb{N}.$ Each $F_n$ is closed, but unbounded (above). Furthermore, for all $n$, $F_{n+1}\subseteq F_n.$ But the intersection $$\bigcap\limits_{n=1}^{\infty} F_n = [1,\infty)\cap[2,\infty)\cap[3,\infty)\cap\ldots$$ is empty, by the following logic. Suppose there is some finite real number $x$ that is in each of the $F_n$. Let $N\in\mathbb{N}$ be the smallest natural number bigger than $x$. Then $x\not\in F_N = [N,\infty)$. But $F_N$ is part of the countably infinite intersection above. So, in fact, $x\not\in\bigcap\limits_{n=1}^{\infty}F_n$, and $\bigcap\limits_{n=1}^{\infty}F_n = \emptyset$.

\vspace{3mm}

\noindent \textbf{(b)} For the second case, consider the sets $$F_n = \left(0, \frac{1}{n}\right),$$ where $n\in\mathbb{N}.$ Each set is open and bounded (above by $1$, and below by $0$). Furthermore, for all $n$, $F_{n+1}\subseteq F_n$. But the intersection $$\bigcap\limits_{n=1}^{\infty} F_n = (0, 1)\cap(0,1/2)\cap(0, 1/3)\cap\ldots$$ is empty. Take any candidate $x$ such that $x\in\bigcap\limits_{n=1}^{\infty}F_n$. let $N\in\mathbb{N}$ be the largest natural number such that $\frac{1}{N} < x$. The set $(0,1/N)$ is part of the countably infinite intersection above, but obviously $x\not\in(0,1/N),$ so $\bigcap\limits_{n=1}^{\infty}F_n = \emptyset.$

\pagebreak

\noindent \textbf{(5.)}

\vspace{3mm}

\noindent\mybox{\colorbox{lightgray}{\textbf{Theorem}}}\hspace{3mm}\hrulefill

\noindent Let $E$ be a set of isolated points. Then $E$ is countable.

\vspace{5mm}

\noindent\mybox{\colorbox{lightgray}{\textbf{Proof}}}\hspace{3mm}\hrulefill

\noindent If $E$ is finite, then $E$ is obviously countable, so suppose $E$ is infinite. Since $E$ consists only of isolated points, each $x\in E$ has an associated radius $r_x>0$ such that $$(x-r_x,x+r_x)\cap E = \lbrace x \rbrace.$$ Note that we can make the $r_x$ sufficiently small so that none of the sets $(x-r_x, x+r_x)$ overlap each other. By making the $r_x$ smaller we do not run the risk of including other points of $E$ (besides $x$) in the interval $(x-r_x, x+r_x)$. Now, with the $r_x$ defined so that no overlapping occurs, for each $x\in E$ define $q_x$ to be some rational number in $(x-r_x, x+r_x)$. The existence of each $q_x$ is guaranteed by the fact that $\mathbb{Q}$ is dense in $\mathbb{R}$. Let $S$ be the set of all the $q_x$. By Theorem 3 from Royden and Fitzpatrick, $S$ is countable, since $S\subseteq\mathbb{Q}$, which is known to be countable. Let $f:E\longrightarrow S$ be defined by $f(x) = q_x$. Clearly, $f$ is onto, since each $q_x\in S$ is defined in terms of an $x\in E$. To show that $f$ is one-to-one, suppose $f(x)=f(y)$, where $x,y\in E$. Then $q_x=q_y$. We know that, if $x\neq y$, then $(x-r_x,x+r_x)$ and $(y-r_y, y+r_y)$ do not overlap, meaning it is impossible for $q_x$ and $q_y$ to be equal (since they are in two different, non-overlapping intervals). Thus, it must be the case that $x=y$, so $f$ is one-to-one. Since $f$ is one-to-one and onto, it is a bijection. Therefore, $E$ is countable.   

\hfill$\blacksquare$

\end{document}


